\documentclass[conference]{IEEEtran}
\IEEEoverridecommandlockouts
\usepackage{amsmath,amssymb,amsfonts}
\usepackage{algorithmic}
\usepackage{graphicx}
\usepackage{textcomp}
\usepackage{xcolor}
\usepackage{booktabs}
\usepackage{cite}
\usepackage{adjustbox}
\usepackage{geometry}
\geometry{a4paper, margin=1in}

% Configuring font package last
\usepackage{times}
\usepackage{url}

\title{Cover Song Identification: Phase 4 Interim Report}

\author{
    \IEEEauthorblockN{Rojin Raju \\ Guided by: Dr. Joe Cheri Ross}
    \IEEEauthorblockA{\textit{National Institute of Technology, Calicut}}
}

\begin{document}

\maketitle

\begin{abstract}
Cover song identification (CSI) aims to recognize alternative renditions of an original song despite variations in timbre, tempo, structure, key, and lyrics. While state-of-the-art models perform well on Western music, their applicability to Indian music—characterized by distinct vocal ornamentations and instrumentation—remains largely unexplored. Following the data collection pipeline established in Phase 3, this Phase 4 report evaluates two prominent deep learning models, ByteCover and TPPNet, on newly validated Indian datasets encompassing Malayalam and Hindi tracks. Through extensive experimentation, including pre-trained benchmarking, model fine-tuning, data augmentation evaluation, and vocal separation testing, we demonstrate that Western-trained models suffer significant performance degradation on Indian music. However, fine-tuning networks explicitly on regional data substantially restores precision. 
\end{abstract}


\section{Introduction}
Cover song identification (CSI) plays a pivotal role in music information retrieval by enabling the detection of different versions of a song, even when altered by changes in key, tempo, instrumentation, or structural composition. Significant strides have been made in optimizing CSI techniques for Western music, leveraging standardized datasets and methodologies. However, extending these methods to Indian music introduces distinct challenges due to its rich and diverse musical heritage. Indian music is characterized by ragas, which are intricate melodic frameworks, and talas, which are complex rhythmic cycles, often accompanied by improvisation and micro-tonal nuances. This report evaluates the performance of CSI models on Indian music and details the development of a comprehensive data collection framework designed to curate a high-quality Indian Cover Song dataset.

\section{Problem Statement}
This study focuses on evaluating the applicability of established cover song identification (CSI) methods to Indian cover songs. Early experiments with models like ByteCover and TPPNet on the current Indian covers dataset revealed that the dataset size and quality were insufficient to draw trustworthy inferences regarding model performance on Indian music. The primary problem addressed in this report is the scarcity of a large-scale, validated Indian cover song dataset. Furthermore, the manual validation of thousands of potential covers is prohibitively time-consuming, necessitating the development of automated validation frameworks that can distinguish between valid covers, exact re-uploads, and unrelated tracks.

\section{Literature Survey}
Our research builds upon two seminal works in cover song identification (CSI) that provide the foundation for evaluating ByteCover and TPPNet. The first, "BYTECOVER: Cover Song Identification via Multi-Loss Training" by Du et al. (2021), introduces ByteCover, a sophisticated model leveraging a ResNet architecture enhanced with Instance-Batch Normalization (IBN) blocks. This design enables the model to extract features resilient to variations in key, tempo, and timbre. ByteCover processes Constant-Q Transform (CQT) spectrograms and utilizes a multi-loss training approach.
The second paper, "Temporal Pyramid Pooling Convolutional Neural Network for Cover Song Identification" by Xu et al. (2019), presents TPPNet, a convolutional neural network employing Temporal Pyramid Pooling (TPP) to capture multi-scale temporal dependencies within CQT spectrograms. TPPNet excels at addressing issues like variable song durations and temporal misalignments.

\section{Datasets}
To comprehensively evaluate the models, we utilized established Western benchmarks alongside our newly compiled Indian datasets.
\begin{itemize}
    \item \textbf{Western Baseline:} SHS100K-TEST, consisting of 2,417 queries.
    \item \textbf{Indian Cover Datasets:}
    \begin{itemize}
        \item \textbf{CoverIndian:} 3,100 tracks (includes augmented versions).
        \item \textbf{Mal (New):} 1,476 tracks of Malayalam covers collected in Phase 3.
        \item \textbf{Hin (New):} 842 tracks of Hindi covers.
        \item \textbf{All Indian:} A combined set totaling 5,418 tracks used to represent the broader Indian musical context.
    \end{itemize}
\end{itemize}



\section{Experiments and Results}
We conducted five targeted experiments to observe how Indian cover songs compare to Western covers under different model configurations. The Ultimate aim is to make sure that any decrease in metrics for Indian Cover songs is purely related to the Indian song features and is independent of other variations in the test data.

\subsection{Exp 1: Pre-Trained Baseline Results}
Both TPPNet and ByteCover, pre-trained strictly on Western songs, were evaluated on our datasets to establish a baseline.

\begin{table}[htbp]
    \caption{Pre-Trained Baseline Results (TPPNet)}
    \centering
    \begin{adjustbox}{width=\columnwidth}
    \begin{tabular}{lccc}
    \toprule
    \textbf{Dataset} & \textbf{MAP} & \textbf{Mean Rank} & \textbf{Total tracks} \\
    \midrule
    SHS100K-TEST & 0.4903 & 11.2 & 2417 \\
    CoverIndian  & 0.3234 & 1.9 & 3100 \\
    Mal          & 0.1757 & 86.8 & 1476 \\
    Hin          & 0.2424 & 14.6 & 842 \\
    All Indian   & 0.2405 & 85.7 & 5418 \\
    \bottomrule
    \end{tabular}
    \end{adjustbox}
\end{table}

\begin{table}[htbp]
    \caption{Pre-Trained Baseline Results (ByteCover)}
    \centering
    \begin{adjustbox}{width=\columnwidth}
    \begin{tabular}{lccc}
    \toprule
    \textbf{Dataset} & \textbf{MAP} & \textbf{Mean Rank} & \textbf{Total tracks} \\
    \midrule
    SHS100K-TEST & 0.0183 & 113.9 & 2417 \\
    CoverIndian  & 0.0052 & 306.5 & 3100 \\
    Mal          & 0.0063 & 455.8 & 1476 \\
    All Indian   & 0.0027 & 852.8 & 5418 \\
    \bottomrule
    \end{tabular}
    \end{adjustbox}
\end{table}
 TPPNet generalizes slightly better to Indian datasets than ByteCover, but still exhibits a significant drop in Mean Average Precision (MAP) compared to Western songs. ByteCover completely fails to generalize, predicting at near-random chance.

\subsection{Exp 2: Fine-Tuned TPPNet Results}
TPPNet was fine-tuned on the "All Indian" dataset using a 60-20-20 split to adapt its feature extraction to Indian musical dimensions.

\begin{table}[htbp]
    \caption{Fine-Tuned TPPNet Results}
    \centering
    \begin{adjustbox}{width=\columnwidth}
    \begin{tabular}{lccc}
    \toprule
    \textbf{Dataset} & \textbf{MAP} & \textbf{Mean Rank}  \\
    \midrule
    All Indian (Pre-Trained) & 0.2405 & 85.7  \\
    \textbf{All Indian (Fine-Tuned)} & \textbf{0.3928} & \textbf{13.60}  \\
    \bottomrule
    \end{tabular}
    \end{adjustbox}
\end{table}
 Explicit training on Indian cover data significantly improves the model's precision and lowers the mean rank error.

\subsection{Exp 3: Data Augmentation Impact (TPPNet)}
TPPNet was evaluated on a unified dataset (7,338 tracks) augmented with tempo and pitch shifts.

\begin{table}[htbp]
    \caption{Data Augmentation Impact (TPPNet)}
    \centering
    \begin{adjustbox}{width=\columnwidth}
    \begin{tabular}{lccc}
    \toprule
    \textbf{Model Origin} & \textbf{MAP} & \textbf{Mean Rank} &  \\
    \midrule
    Western Baseline & 0.4563 & 9.23\\
    \textbf{Indian Fine-Tune} & \textbf{0.4724} & \textbf{8.23}  \\
    \bottomrule
    \end{tabular}
    \end{adjustbox}
\end{table}
Fine-tuning on a fully augmented dataset yields the highest overall performance, from this we can infer that data augmentation is might be yielding a higher result due to the presence of easy to identify queries. Therefore in real world scenarios this might not yield a comparable result.

\subsection{Exp 4: Standalone Clean Dataset Filtering}
Synthesized augmented data was removed to test the models strictly on 4,160 distinct, real-world cover songs.

\begin{table}[htbp]
    \caption{Standalone Clean Dataset Filtering}
    \centering
    \begin{adjustbox}{width=\columnwidth}
    \begin{tabular}{llcc}
    \toprule
    \textbf{Model Checkpoint} & \textbf{MAP} & \textbf{Mean Rank} \\
    \midrule
    TPPNet (Western Baseline)  & 0.3542 & 25.74 \\
    \textbf{TPPNet (Indian Fine-Tune)} & \textbf{0.3704} & \textbf{17.53} \\
    ByteCover (Western Baseline) & 0.0116 & N/A \\
    \bottomrule
    \end{tabular}
    \end{adjustbox}
    \vspace{0.1cm}
    % \scriptsize{*Evaluation Set: Full Clean Indian}
\end{table}
 Performance drops slightly when artificial variations are removed, reflecting the inherent difficulty of real-world acoustic variance.

\subsection{Exp 5: Vocal Separation Performance}
To remove instrumental bias, tracks were stripped to isolated vocals using Demucs, and models were evaluated exclusively on these vocal stems.

\begin{table}[htbp]
    \caption{Vocal Separation Performance}
    \centering
    \begin{adjustbox}{width=\columnwidth}
    \begin{tabular}{lccc}
    \toprule
    \textbf{Dataset} & \textbf{TPP MAP} & \textbf{TPP Mean Rank}  \\
    \midrule
    SHS100K-VOCALS & 0.1761 & 35.61  \\
    INDIAN\_Vocals & \textbf{0.3176} & \textbf{71.59} \\
    \bottomrule
    \end{tabular}
    \end{adjustbox}
\end{table}
 TPPNet relies heavily on instrumentals for Western songs but intrinsically matches Indian vocals better, likely due to distinct, consistent vocal ornamentations present across regional renditions.
 Thus even after the decrease in precision for both datasets, the indian covers where able to score a bit higher than vocally separated western model.

\section{Conclusion and Next Steps}
Successfully bench-marked a newly validated Indian cover song dataset. We confirmed that Western-trained models suffer marked performance drops on Indian music, though fine-tuning networks on regional data substantially restores precision. But overall under most scenarios western covers perform much better than Indian covers. The only test which resulted in a higher accuracy for Indian Covers compared to western covers is when the vocals are separated.

The next steps involve:
\begin{enumerate}
    \item Investigating the specific audio features driving the favorable vocal separation results in Indian tracks.
    \item Developing architectural and methodological improvements tailored specifically to the unique musical dimensions of Indian cover songs.
\end{enumerate}

\begin{thebibliography}{3}
\bibitem{du2021bytecover}
Du, X., et al., ``ByteCover: Cover Song Identification via Multi-Loss Training,'' \textit{Proc. IEEE Int. Conf. Acoust., Speech, Signal Process. (ICASSP)}, 2021.

\bibitem{xu2019tppnet}
Yu, Z., et al., ``Temporal Pyramid Pooling Convolutional Neural Network for Cover Song Identification,'' \textit{Proc. Int. Joint Conf. Artif. Intell. (IJCAI)}, 2019.

\bibitem{yt_dlp_repo}
The yt-dlp developers, ``yt-dlp: A youtube-dl fork with additional features and fixes,'' \textit{[Online]. Available: \url{https://github.com/yt-dlp/yt-dlp}}.
\end{thebibliography}

\end{document}
